\section{帰無仮説検定の応用と注意点}
\label{b14_NHST2}

\subsection{授業内容}
\subsubsection{概要}
前回までに帰無仮説検定の基本的な考え方と手順について学んだ。本コマでは，より実践的な検定法として母分散が未知の場合の平均値の検定(t検定)を学び，検定結果の報告方法や解釈について理解を深める。また，帰無仮説検定における2種類の誤り(第1種の誤りと第2種の誤り)について学び，検定力と効果量の概念を理解する。さらに，帰無仮説検定の限界や注意点についても考察する。特に，統計的有意性と実質的重要性の違い，検定の多重性の問題などについて理解し，帰無仮説検定を適切に活用するための視点を身につける。

\subsubsection{コマ主題細目}
\begin{description}
	\item[1つの平均値の検定(t検定)] 前回は相関係数の検定とカイ二乗検定について学んだが，本コマでは平均値の検定について学ぶ。実際の研究では母分散がわかっていることはほとんどなく，標本から不偏分散を用いて推定する必要がある。このような場合は標準正規分布ではなく，$t$分布と呼ばれる統計量が用いられる。$t$分布の特徴(自由度によって形状が変わる，サンプルサイズが大きくなると標準正規分布に近づくなど)，$t$検定の手順，および$t$分布表の使い方について学ぶ。また，1サンプルの$t$検定と，次回学ぶ2サンプルの$t$検定の違いについても理解する。
	      \begin{flushright}
		      $\to$ \textcite{Kawabata201408}のP.93--96，\textcite{Yamada2004}のP.128--131.
	      \end{flushright}

	\item[検定における2種類の誤り] 帰無仮説検定は確率的な判断を行う手法であるため，誤った判断を下す可能性が常に存在する。検定における誤りには2種類あり，第1種の誤り(Type I Error)は帰無仮説が正しいのに棄却してしまう誤り，第2種の誤り(Type II Error)は帰無仮説が誤りなのに棄却できない誤りである。有意水準αは第1種の誤りを犯す確率の上限を表し，検定力(1-β)は第2種の誤りを犯さない確率を表す。検定力はサンプルサイズ，効果量，有意水準によって決まることを理解し，適切なサンプルサイズの決定方法(検出力分析)についても学ぶ。
	      \begin{flushright}
		      $\to$ \textcite{Yamada2004}のP.120--121，検定の多重性については同書P.158--161，\textcite{Kawabata201408}の P.90--93も参照。
	      \end{flushright}

	\item[効果量と実質的重要性] 統計的有意性(p値が有意水準より小さいこと)は必ずしも実質的な重要性を意味するわけではない。特に大きなサンプルサイズでは，わずかな効果でも統計的に有意になりやすい。そこで，効果の大きさを標準化して表す効果量の概念が重要になる。代表的な効果量として，相関係数(r)，Cohen's d，η²(イータ二乗)などがあり，それぞれの意味と解釈について学ぶ。また，効果量の大きさを評価する目安(小・中・大)についても理解する。効果量は検定結果とともに報告することで，結果の実質的な意味を評価する助けとなる。
	      \begin{flushright}
		      $\to$ 効果量については\textcite{Matia201201}のP.43--46，\textcite{Kawabata201408}のP.97--98.
	      \end{flushright}

	\item[結果の報告と解釈] 帰無仮説検定の結果を論文やレポートで報告する際の適切な方法について学ぶ。統計量の値，自由度，p値，効果量などを正確に報告することの重要性を理解する。また，p値の捉え方として「帰無仮説が正しいという仮定のもとで，観測されたデータやそれよりも極端なデータが得られる確率」であることを再確認し，p値の誤った解釈(帰無仮説が正しい確率，効果の大きさなど)を避けるようにする。さらに，統計的に有意であることの意味と限界について理解を深める。
	      \begin{flushright}
		      $\to$ \textcite{Yamada2004}のP.132--133，\textcite{Yamauchi201001}のP.216--220.
	      \end{flushright}

	\item[帰無仮説検定の限界と注意点] 帰無仮説検定は有用なツールであるが，いくつかの限界や注意点がある。検定の多重性の問題(多数の検定を行うとType I Errorの確率が増加する)，出版バイアス(有意な結果のみが報告される傾向)，p-hacking(有意な結果が得られるまで分析を変更する)などの問題について理解する。また，これらの問題に対する対処法(Bonferroni補正，FDR制御，事前登録など)についても学ぶ。さらに，近年の再現性危機と統計改革の動きについても触れ，帰無仮説検定を補完する方法(信頼区間，ベイズ統計，メタ分析など)についても概観する。
	      \begin{flushright}
		      $\to$ 検定の多重性については\textcite{Yamada2004}のP.158--161，帰無仮説検定の限界については\textcite{Kawabata201408}のP.99--103.
	      \end{flushright}

\end{description}
\subsubsection{キーワード}
\begin{itemize}
	\item t検定と自由度
	\item 第1種の誤りと第2種の誤り
	\item 検定力と効果量
	\item 統計的有意性と実質的重要性
	\item 検定の多重性と帰無仮説検定の限界
\end{itemize}
\subsection{授業情報}
\paragraph{コマの展開方法}
適宜投影資料を用いる
\paragraph{標準シラバスにおける位置づけ}
科目番号4；心理学研究法；2.データを用いた実証的な思考方法；C データの統計的記述\\
科目番号5；心理学統計法；1.心理学で用いられる統計手法；F 推測統計:その考え方
\subsubsection{予習・復習課題}
\paragraph{予習}
前回学んだ帰無仮説検定の基本的な考え方や手順，有意水準とp値の関係について復習しておくこと。特に，相関係数の検定やカイ二乗検定の手順と解釈について確認しておくと理解が深まる。
\paragraph{復習}
授業で学んだ内容を定着させるために，以下の点について理解を深めておくこと：

\begin{enumerate}
\item t検定の手順と適用条件，およびt分布の特徴
\item 第1種の誤りと第2種の誤りの違い，および検定力との関係
\item 効果量の種類と解釈，および統計的有意性との違い
\item 帰無仮説検定の結果を適切に報告する方法
\item 帰無仮説検定の限界と注意点，および対処法
\end{enumerate}

検定の話に入ると，要因計画や線形モデルとどのようにこれらが接合するのかが分かりにくい。回帰モデルの説明変数を離散化したのが要因計画→要因計画の効果の大きさ(傾きの大きさ)を質的に判断する手続きが，平均値差の検定である，という関係であることを整理しておこう。推測統計学は，基本的に母数の推定(回帰係数のパラメータも母数である)をしているのであり，推定量を何らかの基準で判断するのが検定，という繋がりである。もちろん検定を行わずに，大きさをそのまま報告しても良い。むしろ昨今ではそちらの方が推奨されており，些細な違いや「差がないこと」を「差があること」の根拠とすることの問題点が数多く指摘されている。帰無仮説検定の手続きは，過去の論文に示された結果を読むための知識として重要であり，本質である効果の大きさを推定していることを忘れないように。